\begin{afterword}
\summaryhead

The post answers the popular view that rationality opposes emotion. It suggests that people who hold this view are really contrasting fast perceptual judgment with slow deliberation, and it notes that both can be right or wrong. The author's own test is that an emotion is not rational if it rests on mistaken beliefs, and is rational if correct beliefs evoke it. On this test calm has no special standing: fear of a cool iron is a mistake, and so is calm before a hot one.

Emotions, the post says, arise from beliefs. Joy at dreaming that a dead brother is alive turns to sadness on waking, and the sadness is rational. The standard of rationality spreads from a belief (a goblin tied your laces) to the feeling it causes (anger at the goblin). Becoming more rational can make feelings stronger, if you had been denying facts to avoid them. The author recalls long uncertainty about whether strong feeling was allowed, describes a culture that finds it embarrassing, and concludes: ``it is rational to feel.''

\responsehead

The post's central move is a definition, and it is presented as one (``I label''). A feeling is rational when true beliefs evoke it. This removes the assumption that calm is the rational default, which is the post's most useful result. It also settles the conclusion in advance. If a feeling is rational when its beliefs are true, it follows at once that feeling is rational whenever they are. The last paragraph's ``I know, now, that it is rational to feel'' restates the definition.

The test judges only the belief under a feeling. The post grants that anger at the goblin is ``not just about how-the-world-is,'' and stops there. Whether a feeling fits its cause, in kind or in size, is left unjudged. By the test as written, any feeling that a true belief evokes passes, however large.

The account of how feelings work is asserted rather than shown. ``Our emotions arise from our models of reality'' comes with no evidence, and the one example is the easy case: a dream ends, and the feeling changes at once. Feelings that stay after a belief has changed are not discussed, although the last paragraph reports a feeling that persists ``Despite all my philosophy.''

The history is loose. The claim that since Socrates the cultured have hidden strong feeling has no source and is overstated: the eighteenth-century fashion for sensibility prized the display of feeling.

In short: a definition that removes calm's claim to be the rational default, from which the post's conclusion follows directly, and which judges only the belief under a feeling, not whether the feeling fits.
\end{afterword}
